\documentclass{article}
\usepackage[utf8]{inputenc}
\usepackage{hyperref}
\usepackage{listings}
\usepackage{xcolor}
\usepackage{enumitem}

\lstset{
    basicstyle=\ttfamily\small,
    breaklines=true,
    frame=single,
    backgroundcolor=\color{gray!5},
}

\title{Create multiple users from a CSV file containing username and password}
\author{marcos}
\date{2026-04-25}

\begin{document}

\section*{Create multiple users from a CSV file containing username and password}
\textbf{Author:} marcos \hfill \textbf{Date:} 2026-04-25

\noindent Creates multiple user accounts on a Linux system by reading username and password pairs from a CSV file passed as an argument.

\section*{Language}
bash

\section*{Category}
system-administration

\section*{Command}
\begin{lstlisting}
sudo awk -F, 'NR>1 {print $1,$2}' "$1" | while read user pass; do sudo useradd -m "$user" && echo "$user:$pass" | sudo chpasswd; done
\end{lstlisting}

\section*{Explanation}
The command reads a CSV file (passed as first argument) where each row contains a username and password. It skips the header row (NR>1), then for each subsequent row extracts the username and password fields. For each user, it creates a system account with a home directory (-m) and then uses chpasswd to set the plain text password from the CSV file.

\section*{Tags}
user-management, bulk-operations, system-administration, csv, automation

\section*{Dependencies}
coreutils, shadow-utils, openssl


\section*{Arguments}
\begin{enumerate}
  \item \textbf{CSV\_FILE} (Required): Path to CSV file with username,password format
\end{enumerate}

\section*{Examples}
\begin{enumerate}
  \item \texttt{echo -e 'username,password\textbackslash\{\}nuser1,pass123\textbackslash\{\}nuser2,pass456' > users.csv \&\& sudo awk -F, 'NR>1 \{print \$1,\$2\}' "users.csv" | while read user pass; do sudo useradd -m -p \$(openssl passwd -1 "\$pass") "\$user"; done} - Create users from users.csv file with username,password format
  \item \texttt{sudo awk -F, 'NR>1 \{print \$1,\$2\}' "/path/to/users.csv" | while read user pass; do sudo useradd -m -p \$(openssl passwd -1 "\$pass") "\$user"; done} - Create users from a specific CSV file path
  \item \texttt{alias bulkaddusers='sudo awk -F, '"'"'NR>1 \{print \$1,\$2\}'"'"' "\$1" | while read user pass; do sudo useradd -m -p \$(openssl passwd -1 "\$pass") "\$user"; done'} - Create an alias for bulk user creation
\end{enumerate}

\section*{Output}
\begin{lstlisting}
Creating users from users.csv:
user1: user created
user2: user created
\end{lstlisting}

\section*{Notes}
\begin{itemize}
  \item The CSV file should have a header row and contain username,password in each subsequent row
  \item Requires sudo privileges to create users
  \item Uses openssl to properly hash passwords for useradd
  \item The -m flag creates a home directory for each user
  \item The command skips the first line (header) with NR>1
  \item For better security, consider using a more robust password hashing method
  \item The CSV file must be passed as the first argument to the command
\end{itemize}

\section*{Warnings}
\begin{itemize}
  \item This command requires sudo access and will modify system user accounts
  \item Passwords are passed in cleartext in the CSV file which is a security risk
  \item The default password hashing method (openssl passwd -1) uses MD5 which is not the most secure
  \item Always backup your system before running bulk user creation
  \item Ensure CSV file permissions are secured to prevent unauthorized access
  \item Test with a single entry first before running on multiple users
  \item Consider using a more secure method for password handling in production environments
\end{itemize}

\section*{See Also}
\begin{itemize}
  \item user-activity-and-quota-report
  \item nonhuman-user-activity-and-quota-report
\end{itemize}

\section*{Status}
draft

\section*{Safety}
caution

\section*{Shell}
bash

\section*{Platforms}
linux, gnu-linux

\section*{Created At}
2026-04-25

\section*{Updated At}
2026-04-25

\end{document}
